% !TEX root = tarasov-qasp-zh.tex

\chapter[你怎样解释物体的失重？]{你怎样解释物体的失重？\label{ch-09}}

\begin{dialogue}

\tchr 你对下面这个说法是怎么理解的：\hlred{在行星的赤道上，物体的重量比在两极的小？}

\stdb 我是这么理解的。地球对物体的吸引力在赤道上比在两极小，这有两个原因。第一，地球在两极处稍扁，因此从地球中心到两极的距离比到赤道的距离稍小。第二，地球绕自转轴旋转，其结果是赤道处的吸引力因离心效应而减弱。

\stda 请把你最后那句话讲得再清楚些。

\stdb 这就是：必须从吸引力中减去离心力（centrifugal force）。

\stda 我不同意。第一，\hlblue{离心力并不是作用在做圆周运动的物体上的力。}这一点我们在上一章（\chap{ch-08}）已经讨论过了。第二，即使存在这样的力，它也不可能使吸引力变得与地球不转动时有所不同。吸引力等于 $GmM/r^{2}$，就其本身而言，它并不因为物体还可能受到别的力的作用而改变。

\tchr 你们看，“物体的有重性”（weightness）问题并不像初看起来那么简单。正因如此，它属于考生经常答错的问题之列。事实上，如果我们约定“物体的重量”（weight）是指地球吸引物体的力，即力 $GmM/r^{2}$，那么赤道上重量的减小就只应归因于两极的扁缩（或者说赤道的隆起）。
\pnote{原书在此自造了 ``weightness'' 一词（“失重”weightlessness 的反义构词），指物体处于“有重量”的状态；中译作“有重性”，下文“随意处置物体的‘有重性’”即随意处置其重量。}

\stdb 可是，地球的自转不能置之不理啊！

\tchr 我完全同意。不过我想先指出：通常\hlblue{在日常生活中，“物体的重量”不是被理解为地球吸引它的力——那样理解是完全合乎逻辑的——而是被理解为弹簧秤测出的力，即物体压在地球上的力。}换句话说，测出的是支承力（bearing reaction）（按牛顿第三定律，物体压在支承物上的力与支承力大小相等）。由此可见，“物体在赤道上比在两极轻”这个说法的意思是：物体在赤道上压在支承物上的力比在两极的小。
\pnote{本章反复区分两种“重量”：地球的吸引力，即万有引力 $F=GMm/r^{2}$（人教版必修第二册第七章《万有引力与宇宙航行》）；与弹簧秤测出的“压在支承物上的力”，即视重。人教版必修第一册第三章把重力定义为“由于地球的吸引而使物体受到的力”，第四章第 6 节“超重和失重”讨论的“物体对支持物的压力（或对悬挂物的拉力）”正是第二种——超重、失重说的都是它的增减。量级上：地面上赤道与两极的重力加速度差约 $0.05~\mathrm{m/s^{2}}$（$9.78$ 与 $9.83~\mathrm{m/s^{2}}$），其中自转的贡献约 $0.03~\mathrm{m/s^{2}}$，两极扁缩的贡献约 $0.02~\mathrm{m/s^{2}}$。}

我们把两极处的吸引力记作 $P_{1}$，赤道处的吸引力记作 $P_{2}$；两极处的支承力记作 $N_{1}$，赤道处的支承力记作 $N_{2}$。在两极，物体静止；在赤道，物体做圆周运动。于是
\begin{equation*}%
\begin{split}
P_{1} - N_{1} & =0\\
P_{2} - N_{2}& =ma_{cp}
\end{split}
\end{equation*}
其中 $a_{cp}$ 是向心加速度（centripetal acceleration）。可以把这些方程改写成
\begin{equation}%
\left.
\begin{split}
N_{1} & = P_{1} \\
N_{2} & = P_{2} - ma_{cp}
\label{eq-46}
\end{split}
\right\}
\end{equation}
这里可以清楚地看出 $N_{2}$ 小于 $N_{1}$：第一，$P_{2}$ 小于 $P_{1}$（两极扁缩的效应）；第二，还要从 $P_{2}$ 中减去一个量 $ma_{cp}$（地球自转的效应）。
\pnote{底本此处把赤道处的支承力也误记作 $P_{2}$，按下面的方程组应为 $N_{2}$，中文版径改。另外，本章开篇教师所引的话底本作 ``a body weightless than at the poles''，显系漏排动词（应为 ``a body weighs less than at the poles''），中文版亦径改。}

\stdb 这么说，“物体失去了一半重量”并不意味着地球（或任何其他行星）吸引它的力减小了一半？

\tchr 是的，并不意味着。吸引力可以完全没有改变。这个说法的意思是：\hlblue{物体压在支承物上的力（换句话说，支承力）减小了一半。}

\stdb 可是这样一来，我就可以随心所欲地处置物体的“有重性”了。有什么能阻止我在物体下面挖一口深井，让物体连同它的支承物一起落进井里去呢？这时，将没有任何力压在支承物上。这是否意味着物体“完全失去了重量”？意味着它处于失重（weightlessness）状态？

\tchr 你自己独立地得出了正确的结论。事实上，\hlblue{失重状态就是物体下落的状态。}关于这一点，我想说几句。我见到过把失重解释成“地球的引力被另外某个力平衡的状态”：就卫星的情形而言，充当这个平衡力的说的是离心力。这种说法讲的是：地球吸引卫星的力与离心力互相平衡，结果作用在卫星上的合力等于零，而这正对应于失重。

你们现在当然明白，这种讲法是不对的，单凭卫星并没有受到离心力这一点就足够了。顺便说一句，\hlblue{如果把失重理解为引力被另外某个力平衡的状态，那么物体干脆静止在水平面上时，说它失重倒更有道理。}这恰恰正是重量被支承力平衡的一种情形！实际上，失重根本不需要引力的平衡。恰恰相反，\hlblue{要使物体失重，必须创造这样的条件：除引力外，没有任何别的力作用在物体上。}换句话说，必须使支承力等于零。物体在引力作用下的运动就是下落。因此，失重是一种下落状态，例如矿井中升降机的下落，或者卫星绕地球的运行。

\stda 在上一章（\chap{ch-08}）里您说过，卫星绕地球运行，无非是在无限长的时间里落向地球。

\tchr 卫星绕地球的运动是下落，这一点可以用下面的办法非常直观地表明。设想你站在一座高山的山顶上，沿水平方向抛出一块石头。我们将忽略空气阻力。石头的初速度越大，它落得越远。\figr{fig-42}~\drkgry{(a)} 表示了随着初速度增大，石头轨迹逐渐变化的情形。在某个速度 $v_{1}$ 下，下落石头的轨迹变成圆，石头就成了地球的卫星。速度 $v_{1}$ 称为环绕速度（circular orbital velocity）。它由式~\eqref{eq-34} 得出：
\begin{equation}%
%\left.
%\begin{split}
v_{1} = \sqrt{\frac{GM}{r}}
\label{eq-47}
%\end{split}
%\right\}
\end{equation}
\begin{figure}[!h]%
\centering
\includegraphics[width=0.44\textwidth]{figs/ch-09/fig-042a.pdf}
\caption[初速度不同时物体绕地球运动的轨迹。]{初速度不同时物体绕地球运动的轨迹。\label{fig-42}}
\end{figure}
若近似取卫星轨道半径 $r$ 等于地球半径，则 $v_{1} \approx \SI[per-mode=symbol]{8}{\kilo\metre\per\second}$。
\pnote{环绕速度即人教版必修第二册第七章第 4 节“宇宙航行”中的第一宇宙速度，教材给出的数值为 7.9~km/s（原书此处取近似值 8~km/s）。}

\stda 如果从山顶抛石头时继续增大初速度，会发生什么？

\tchr 石头将沿越来越扁长的椭圆绕地球运行（\figr{fig-42}~\drkgry{(b)}）。在某个速度 $v_{2}$ 下，石头的轨迹变成抛物线，石头就不再是地球的卫星了。速度 $v_{2}$ 称为脱离速度（escape velocity）。据计算，$v_{2} \approx \SI[per-mode=symbol]{11}{\kilo\metre\per\second}$，约为 $v_{1}$ 的 $\sqrt{2}$ 倍。
\pnote{脱离速度即人教版同一节（“宇宙航行”）中的第二宇宙速度，数值为 11.2~km/s。注意：人教版把“逃逸速度”一名留给飞出太阳系的第三宇宙速度（16.7~km/s），故本书的 escape velocity 译作“脱离速度”，以免与人教版用语相混。}

\stda 您把失重状态定义为下落的状态。然而，如果石头的初速度达到脱离速度，石头将离开地球。这种情形下不能说它在落向地球。那么，石头的失重该怎么解释呢？

\tchr 很简单。这种情形下的失重，是石头落向太阳。
\pnote{按本章的用法，“落向太阳”应理解为“只在太阳引力作用下的运动”：石头脱离地球后是在绕太阳的椭圆轨道上运行（与地球一样），并不会真的落到太阳上——正如“落向地球”的卫星不会落到地面上。}

\stda 那么，位于星际空间某处的宇宙飞船，它的失重就应归之于飞船在某个天体的引力场（gravitational field）中的下落了？

\tchr 正是。

\stdb 不过我觉得，把失重定义为下落状态还需要加以精确化。跳伞员也在下落，可他丝毫没有失重的感觉。

\tchr 你说得对。失重并不是随便哪一种下落。失重是所谓的自由落体（free fall），即物体仅（！）在引力作用下发生的运动。我已经说过，要使物体失重，必须创造除引力外没有别的力作用在物体上的条件。跳伞员下落时，就有一个附加的力——空气阻力。
\pnote{人教版中“自由落体运动”指物体只受重力、由静止开始下落的运动（必修第一册第二章）；本章取其广义——只受引力作用即可，初速度可以是任意值（水平抛出的石头、卫星都算）。原文“仅（！）……”强调的正是“只受引力”这一条件。}
\pnote{跳伞员离机后、开伞前的一段近似自由下落中确有短暂的近失重感；教师指的是整个下落过程——尤其是开伞后接近匀速下降的阶段——始终受空气阻力作用，因而不是失重。}

\end{dialogue}


\section*{习题}
\label{problems-09}
\addcontentsline{toc}{section}{习题}

\begin{enumerate}[resume*=problems]
\item 已知某球状行星的自转周期为 10 小时，且物体在该行星的赤道上处于失重状态。求该行星物质的密度。
\pnote{题设“赤道上失重”即式~\eqref{eq-46} 中 $N_{2}=0$ 的情形：$GmM/R^{2}=m\omega^{2}R$（$R$ 为行星半径）。由此可得行星密度 $\rho=3\pi/(GT^{2})$，与行星半径无关；代入 $T=10$ 小时，得 $\rho\approx 110~\mathrm{kg/m^{3}}$，比水还小——能在赤道失重的行星必然极其蓬松。}
\end{enumerate}
%
\cleardoublepage

\thispagestyle{empty}
\vspace*{-2.2cm}
\begin{figure*}[!ht]
\centering
\includegraphics[width=0.68\textwidth]{figs/sec/tarasov-qasp-sec-05.pdf}
\end{figure*}
%\begin{centering}
\begin{fullwidth}
%%\paragraph{}
{\textsf{\Large \hlred{守恒定律的作用是怎么估计也不会过高的。它们是人类在许多代人经验的基础上确立起来的最普遍的法则。熟练地运用守恒定律，可以使许多问题的求解变得相当容易。下面我们来考察一些与能量守恒和动量守恒定律有关的例子。
}}}
\end{fullwidth}

% NEWTERMS: weightlessness | 失重 | 章主题；对应人教版必修第一册第四章"超重和失重"
% NEWTERMS: weight | 重量 | 本章区分"重量"两种定义（地球引力／弹簧秤所测"压在支承物上的力"，即视重）；第 2 章力图语境沿用"重力"
% NEWTERMS: weightness | 有重性 | 原书自造词，weightlessness 的反义构词，指"有重量"的状态
% NEWTERMS: centrifugal force | 离心力 | 承第 8 章：不作用于做圆周运动的物体
% NEWTERMS: centripetal acceleration | 向心加速度 | 原书符号 $a_{cp}$
% NEWTERMS: circular orbital velocity | 环绕速度 | 即第一宇宙速度（7.9 km/s）
% NEWTERMS: escape velocity | 脱离速度 | 即第二宇宙速度（11.2 km/s）；人教版"逃逸速度"指第三宇宙速度，故不译"逃逸速度"
% NEWTERMS: free fall | 自由落体 | 本章取广义：只受引力作用的运动
% NEWTERMS: gravitational field | 引力场 | 
% NEWTERMS: force of attraction | （地球的）引力 | 即万有引力 $F=GMm/r^{2}$
